% !TEX root = ../Presentation.tex
% ---------------------------------------------------------------------------
% Model A: sessions pooled
% Slides taken from ../Presentation_ABC/sections; trimmed ones are marked.
% ---------------------------------------------------------------------------
\section{Model A: Sessions Pooled}

\subsection{Model Structure}

\begin{frame}{Model A: graphical model}
  \begin{columns}[c]
    \begin{column}{0.6\textwidth}
      \centering
      \begin{tikzpicture}[gm, node distance=0.7cm and 0.6cm]
        \node[latent] (b0) {$\beta_0$};
        \node[latent, right=of b0] (b1) {$\beta_1$};
        \node[latent, right=of b1] (gh) {$\gamma_h$};
        \node[latent, right=of gh] (gl) {$\gamma_l$};

        \node[det, below=1.0cm of $(b1)!0.5!(gh)$] (psi) {$\psi_i$};
        \node[obs, below=0.6cm of psi] (r) {$r_i$};
        \node[const, right=0.8cm of psi] (x) {$x_i$};

        \foreach \p in {b0, b1, gh, gl}
          \draw[arr] (\p) -- (psi);
        \draw[arr] (x) -- (psi);
        \draw[arr] (psi) -- (r);

        \node[plate, fit=(x)(psi)(r), inner sep=5pt] (ptrial) {};
        \node[font=\scriptsize, anchor=north east] at (ptrial.south east) {trials $i = 1,\dots,n$};
        \node[plate, fit=(b0)(gl)(ptrial), inner sep=14pt] (pgroup) {};
        \node[font=\scriptsize, anchor=south west] at (pgroup.south west) {each of the 4 groups};
      \end{tikzpicture}
    \end{column}
    \begin{column}{0.38\textwidth}
      \footnotesize
      \begin{itemize}
        \item circles: random variables
        \item shaded: observed
        \item rectangles: deterministic
        \item $x_i$: fixed covariate
        \item outer plate: the same model, fit independently to each group
      \end{itemize}
    \end{column}
  \end{columns}
\end{frame}

\subsection{Synthetic Data}
% numbers from Model_A/Convergence_A_Report.tex

% trimmed
\begin{frame}{Model A on synthetic data: what is recoverable}
  At the size of the real data (about 3{,}230 trials per group):
  \begin{itemize}
    \item \textbf{Recoverable:} PSE, JND, $\gamma_h$, unbiased and with calibrated 95\% HDIs
    \item \textbf{Not reliably recoverable:} small $\gamma_l$, biased upward
    \item \textbf{Caveat:} data simulated with one curve per group, so this says nothing about
          variation between sessions
  \end{itemize}
\end{frame}

\subsection{Real Data}
% numbers from Model_A/Fit_A_Report.tex

\begin{frame}{Model A on the real data: fitted curves}
  \centering
  \includegraphics[width=0.48\textwidth]{\plotsA/Fit_Plots_A/curves_left.png}\hfill
  \includegraphics[width=0.48\textwidth]{\plotsA/Fit_Plots_A/curves_right.png}
  \par\vspace{0.4em}{\small Unimanual in blue, bimanual in red, with the observed proportions
  pooled over sessions; dashed: the trained threshold of 28~$\mu m$.}
\end{frame}

% Distractor effect, one slide per quantity: its row of uni_bi_posteriors.png (unimanual and
% bimanual posteriors overlaid) above its row of effect_densities.png (bimanual - unimanual).
% Both PNGs are 2223 px wide at 200 dpi (1 px = 0.36 bp) with the same two hand columns;
% the rows are cropped in LaTeX at the white gaps between them, hand titles left out.
\newcommand{\effectpair}[4]{%
  \begin{frame}{Model A on the real data: effect of the distractor on #1}
    \centering
    \begin{tikzpicture}
      \node[inner sep=0] (ub) {\includegraphics[width=0.84\textwidth, #2]{\plotsA/Fit_Plots_A/uni_bi_posteriors.png}};
      \node[inner sep=0, below=3pt of ub] (ef) {\includegraphics[width=0.84\textwidth, #3]{\plotsA/Fit_Plots_A/effect_densities.png}};
      \node[above, font=\scriptsize] at ($(ub.north west)!0.255!(ub.north east)$) {Left Hand};
      \node[above, font=\scriptsize] at ($(ub.north west)!0.752!(ub.north east)$) {Right Hand};
      \node[left, rotate=90, anchor=south, font=\scriptsize, muted] at (ub.west) {uni vs bi};
      \node[left, rotate=90, anchor=south, font=\scriptsize, muted] at (ef.west) {bi $-$ uni};
    \end{tikzpicture}\\[0.1em]
    {\scriptsize \textcolor{plotblue}{\rule[0.5ex]{1.2em}{1pt}}~Unimanual\quad
     \textcolor{plotred}{\rule[0.5ex]{1.2em}{1pt}}~Bimanual\quad
     {\color{muted}95\% HDI shaded, mean dotted#4}}
  \end{frame}}

\effectpair{PSE}
  {trim=0 626.04bp 0 38.88bp, clip}
  {trim=0 572.76bp 0 38.88bp, clip}
  {; green dashed: trained threshold}

% trimmed
\begin{frame}{Model A: conclusions}
  \begin{itemize}
    \item The pooled data are reproduced at every amplitude
    \item Effect of the distractor (bimanual $-$ unimanual):
    \begin{itemize}
      \item \textbf{PSE:} lowered, by about 4~$\mu m$ (left) and 6~$\mu m$ (right)
      \item \textbf{JND:} raised for the left hand; unclear for the right
      \item \textbf{\boldmath$\gamma_h$:} raised, by about 0.04
      \item \textbf{\boldmath$\gamma_l$:} no clear effect
    \end{itemize}
    \item[$\Rightarrow$] A curve for every session: Model B
  \end{itemize}
\end{frame}
